We built a judge and argued it is the instrument we stake the book on.
Chapter 6 showed how we test it (mutation testing) and that it caught a
bug. This chapter is the aggressive version: we stop politely checking
the judge and start deliberately trying to \emph{break} it --- the way a
determined detractor (or a suspicious reader) would.

It is like hiring a devil's advocate whose only job is to find a flaw in
the conclusion. We want to know how confident we can be in any number
our judge produces --- and the only way to find out is to attack it.

There is a reason to be particularly aggressive here, and it goes to the
heart of the book. Our whole method is built on trusting measurements:
we claim the model is surface-tied, non-compositional, that it forgets
under sequential training. Every one of those claims passes through the
judge. So if the judge is broken in even one respect, every conclusion
in the book could be subtly wrong. Red-teaming the judge is not a single
chapter's task; it is the standing obligation of the whole experiment.
And it is an \emph{unbounded} one --- we can never prove the judge is
perfect, only that we have tried hard to break it and failed.

\section{The attack surface}\label{the-attack-surface}

A judge that executes generated code is attacked through the code it is
given. We enumerate the ways a candidate could slip through:

\begin{itemize}
\tightlist
\item
  \textbf{hardcoded outputs} --- a program that ignores the input and
  returns a fixed answer that happens to match one test,
\item
  \textbf{mutation variants} --- operator flips, wrong variables,
  off-by-one, missing edge-cases (the mutation suite from Chapter 6),
\item
  \textbf{weak test cases} --- if the hidden tests are too few or too
  narrow, a wrong program can pass,
\item
  \textbf{boundary failures} --- behavior at zero, negatives, or large
  values that the tests don't probe,
\item
  \textbf{resource abuse} --- programs that loop forever or consume too
  much,
\item
  \textbf{unsupported syntax} --- smuggling in a construct the sandbox
  should reject,
\item
  \textbf{ambiguous specifications} --- a prompt the model could
  reasonably read two ways,
\item
  \textbf{evaluator inconsistencies} --- the same program scored
  differently in two runs.
\end{itemize}

This list is a map of the ways the judge could be fooled --- and each
entry is a place we must probe. The point of enumerating the surface is
that skepticism is only useful when it is specific. ``The judge might be
wrong'' is not actionable. ``The judge might let a hardcoded-output
program pass if the tests are too narrow'' is a concrete thing we can
check.

\section{What we found}\label{what-we-found}

The mutation suite is the sharpest tool here. By generating
deliberately-broken variants of the reference solution and checking that
the hidden tests catch them, we measured how discriminating the tests
really are. If an obviously-wrong mutant --- say \texttt{+} changed to
\texttt{-} --- passes our tests, the tests are too weak and must be
strengthened until important mutant classes are reliably detected.

A useful way to think about mutation testing is that it measures the
\emph{discriminating power} of the test suite. A test suite that passes
for a correct program and fails for an incorrect one is a good suite. A
suite that passes for both is not testing anything --- it is going
through the motions. The mutation score tells us, concretely, how much
of that ``should be caught'' space our tests actually cover. And the
whole point is to drive it up: the more wrong programs the judge
rejects, the more we trust that a program it accepts is actually right.

And the judge's honesty was validated the hard way: it reported the
model at 0.000 on our very first evaluation, and that was a \emph{judge
bug}, not a model failure. Once we fixed the judge, the model's true
0.897 emerged. A judge that can be wrong --- and we proved ours was ---
must be red-teamed before its numbers are trusted.

\section{The honest framing}\label{the-honest-framing}

Red-teaming the judge is not a formality; it is the difference between
``we ran a number'' and ``we know the number means something.'' The
specific lessons:

\begin{itemize}
\tightlist
\item
  \textbf{Strengthen tests until wrong programs fail.} The mutation
  score is the measure.
\item
  \textbf{Trust the judge less than the tests.} The judge is a pipeline;
  the \emph{hidden tests} are what actually catch errors, so they must
  be numerous and edge-case-complete.
\item
  \textbf{Attack our own instrument deliberately.} The adversarial habit
  --- ``how could this number be wrong?'' --- is what keeps the whole
  experiment honest.
\end{itemize}

A judge that has been red-teamed is one we can believe. A judge we've
only used is one we cannot.

And there is a final lesson that reaches beyond judges. Every
measurement instrument --- in science, in engineering, in daily life ---
can fail in ways its designer did not anticipate. The discipline of
attacking our own apparatus is the same discipline that makes a result
worth believing. It is not paranoia; it is the price of a claim we can
stand behind. The book's whole credibility rests on having paid it, and
having recorded that we paid it, rather than merely asserting that we
are trustworthy.
